\section*{Chapitre \ref{ch:g1:landscapes-around-us} --- Les paysages autour de nous}

\begin{solution}{exo:g1:landscapes-around-us:1}
La \omterm{def:g1:landscapes-around-us:landform}{montagne}, la \omterm{def:g1:landscapes-around-us:landform}{colline}, la \omterm{def:g1:landscapes-around-us:landform}{vallée}, la \omterm{def:g1:landscapes-around-us:landform}{plaine}.
\end{solution}

\begin{solution}{exo:g1:landscapes-around-us:2}
Une \omterm{def:g1:landscapes-around-us:landform}{montagne}~: elle est plus haute et plus raide qu'une \omterm{def:g1:landscapes-around-us:landform}{colline}.
\end{solution}

\begin{solution}{exo:g1:landscapes-around-us:3}
Naturels~: la \omterm{def:g1:landscapes-around-us:water}{rivière}, la forêt, la plage. Construits~: la route, le pont,
la maison.
\end{solution}

\begin{solution}{exo:g1:landscapes-around-us:4}
Là où la terre rencontre la mer.
\end{solution}

\begin{solution}{exo:g1:landscapes-around-us:5}
Par exemple~: naturels, les \omterm{def:g1:landscapes-around-us:landform}{montagnes}, la forêt sur les pentes ou le
ruisseau~; construits, une maison ou l'église du village.
\end{solution}

\begin{solution}{exo:g1:landscapes-around-us:6}
La \omterm{def:g1:landscapes-around-us:landform}{montagne}, avec de la neige à son sommet. La \omterm{def:g1:landscapes-around-us:water}{rivière} coule au fond de la
\omterm{def:g1:landscapes-around-us:landform}{vallée}, entre les \omterm{def:g1:landscapes-around-us:landform}{montagnes} et la \omterm{def:g1:landscapes-around-us:landform}{colline}.
\end{solution}

\begin{solution}{exo:g1:landscapes-around-us:7}
Une \omterm{def:g1:landscapes-around-us:water}{rivière} coule des terres hautes (\omterm{def:g1:landscapes-around-us:landform}{collines} et \omterm{def:g1:landscapes-around-us:landform}{montagnes}) jusqu'à un \omterm{def:g1:landscapes-around-us:water}{lac}
ou jusqu'à la mer. L'eau descend toujours la pente, donc elle ne peut pas
monter jusqu'en haut d'une \omterm{def:g1:landscapes-around-us:landform}{colline}.
\end{solution}

\begin{solution}{exo:g1:landscapes-around-us:8}
L'arbre est près de nous, au premier plan~; la \omterm{def:g1:landscapes-around-us:landform}{colline} est loin, à
l'arrière-plan. Ce qui est loin paraît plus petit que ce qui est proche.
\end{solution}

\begin{solution}{exo:g1:landscapes-around-us:9}
Premier plan~: de hautes herbes et quelques buissons en bas de l'image.
Second plan~: la \omterm{def:g1:landscapes-around-us:water}{rivière} qui serpente à travers la \omterm{def:g1:landscapes-around-us:landform}{plaine}, bordée
d'arbres, des champs, une ferme et une route.
Arrière-plan~: la terre plate qui continue jusqu'à l'horizon, et le ciel.
\end{solution}

\begin{solution}{exo:g1:landscapes-around-us:10}
Parce que ce sont les gens qui l'ont fait~: ils ont défriché la terre,
l'ont labourée et y ont semé le blé. Sans les gens, il y aurait là de
l'herbe, des buissons ou une forêt, et non un champ de blé.
\end{solution}

\begin{solution}{exo:g1:landscapes-around-us:11}
Pas changé~: la \omterm{def:g1:landscapes-around-us:landform}{colline} et la \omterm{def:g1:landscapes-around-us:water}{rivière}. Ajouté par les gens~: les champs,
d'autres maisons, la route et le pont (et la plus grande partie de la
forêt a été coupée).
\end{solution}
